%!TEX program = pdflatex
% Mathematics article -- amsart, the neutral class of the subject. Drafting
% here and converting on acceptance costs nothing, because a journal class
% changes the front matter and the page layout and leaves the body alone.
\documentclass[11pt,reqno]{amsart}

% amsart's own page is set for the AMS printed page and reads narrow on
% screen; a journal class resets it at submission.
\usepackage[margin=1in]{geometry}

% amsart already loads amsmath, amsthm and amsfonts.
\usepackage{amssymb,mathtools}
\usepackage{graphicx}
\usepackage{bm}
\usepackage{xcolor}
% Sorts and compresses each multiple citation, so \cite{c,a,b} prints as
% [1--3] rather than in the order the keys were typed.
\usepackage{cite}
\usepackage{hyperref}
% Links are coloured rather than boxed: the default frames print as
% rectangles around every citation and cross-reference.
\hypersetup{colorlinks=true,
            linkcolor=blue!55!black,
            citecolor=blue!55!black,
            urlcolor=blue!55!black}
\usepackage[capitalize,nameinlink]{cleveref}

% One shared counter, so a Lemma 2.3 and a Theorem 2.4 never collide.
\theoremstyle{plain}
\newtheorem{theorem}{Theorem}[section]
\newtheorem{lemma}[theorem]{Lemma}
\newtheorem{proposition}[theorem]{Proposition}
\newtheorem{corollary}[theorem]{Corollary}
\newtheorem{conjecture}[theorem]{Conjecture}

\theoremstyle{definition}
\newtheorem{definition}[theorem]{Definition}
\newtheorem{example}[theorem]{Example}

\theoremstyle{remark}
\newtheorem{remark}[theorem]{Remark}

% amsart numbers equations straight through; section numbering matches the
% theorem counter and survives inserting a section.
\numberwithin{equation}{section}

\title{Automatic Moment Control for Self-Consistent Oscillator Branches}

\author{Zhiyuan Yao}
\address{Lanzhou Center for Theoretical Physics, Key Laboratory of Theoretical
  Physics of Gansu Province, Key Laboratory of Quantum Theory and Applications
  of MoE, Gansu Provincial Research Center for Basic Disciplines of Quantum
  Physics, Lanzhou University, Lanzhou 730000, China}
\email{yaozy@lzu.edu.cn}

\date{\today}

% AMS journals require a 2020 Mathematics Subject Classification, primary
% first; keywords are optional and printed with it.
% \subjclass[2020]{Primary 00A00; Secondary 00B00}
% \keywords{KEYWORDS}

\begin{document}

% amsart typesets the abstract from the title block, so it is
% declared inside the document and before that block.
\begin{abstract}
For the finite-dimensional self-consistently renormalized oscillator of
Baez and Zhou, local uniqueness was stated under a bounded interaction-moment
hypothesis.  We prove that this hypothesis follows from norm continuity of
the ground-state branch.  State-relative Wick polynomials form a formal
multivariate Appell family.  If a norm-convergent sequence has an escaping
top moment, total-variation concentration makes all lower moments
subhomogeneous, and changing the reference state forces the Wick expectation
in the limit state to tend to $-\infty$.  Ground-state minimality supplies an
opposing uniform lower bound, including at zero coupling through the free
vacuum.  Thus the top moment is bounded on every compact coupling interval.
The cone inverse theorem then yields a uniform neighborhood on which every
shorter closed interval carries exactly one norm-continuous nonnegative
Schwartz ground-state branch, without an a priori moment bound.
\end{abstract}

\maketitle

\section{Introduction}

Ordinary normal ordering originates in Wick's contraction formula
\cite{Wick_1950te}; Gaussian Wick polynomials and constructive field-theory
implementations are treated, from complementary viewpoints, in
\cite{Janson_Book,Glimm_Book,Simon_Book,Simon_1972hs}.  The products in this
paper are different: their reference state need not be Gaussian, and their
centering and commutator axioms make them a state-relative formal Appell
family.

This state-relative construction grew out of Segal's weak-process calculus
and his construction of nonlinear local quantum processes
\cite{Segal_1969nf,Segal_1970nf,Segal_1970co,Segal_1971co}; a broader account
is given by Baez, Segal, and Zhou \cite{Baez_Book}.  Friedman isolated the
corresponding renormalized oscillator equation and its translation
obstruction \cite{Friedman_1973ro}.  Baez and Zhou then formulated
finite-dimensional self-consistent Wick ordering as a nonlinear ground-state
problem and proved a local branch theorem for coercive polynomial
interactions \cite{BaezZhou1992}.

A self-consistent ground-state equation, in which the operator depends on the
state it is meant to produce, is the setting of Hartree-type problems, and the
uniqueness of such ground states has its own literature, and in it the
answers are conditional: Lenzmann settled the pseudo-relativistic Hartree
equation for ground states of sufficiently small mass, and there for all but
countably many values of it \cite{Lenzmann_2009uo}, while Griesemer and
Hantsch settled the Hartree--Fock ground state of a closed-shell atom whose
nuclear charge is large enough relative to its electron number
\cite{Griesemer_2011us}; the question remains active for nonlinear Hartree
equations \cite{Luo_2018uo}.  What distinguishes the present problem
from all of these is that the nonlinearity enters through the Wick ordering
rather than through a convolution, so the interaction moment that has to be
controlled is the one the ordering itself produces.

The analytic-family and self-adjoint Schr\"odinger background is supplied by
standard perturbation and operator theory
\cite{Kato_Book,ReedSimonII,Simon_1970cc}.  Baez and Zhou's uniqueness
argument assumes that the interaction moment remains bounded along the
branch, and they record the status of that assumption themselves: they call
it a technical condition and ask whether it can be omitted, while noting that
continuity on the closed interval cannot be, since uniqueness already fails
on a half-open one \cite{BaezZhou1992}.  This paper answers that
question.  Since norm convergence alone does not normally control unbounded
polynomial observables, the answer uses the self-consistent variational
structure rather than a compactness slogan.

The mechanism is a sign contradiction.  Formal Appell calculus identifies
the leading behavior of a Wick polynomial when its reference state develops
an escaping top moment.  Total-variation convergence of position laws makes
all lower-moment contributions negligible, leaving a negative coercive term.
The same expectation, evaluated in a fixed comparison state, is bounded
below by ground-state minimality.  This also handles sequences of couplings
approaching zero, where the free vacuum is the comparison state.

\begin{theorem}\label{thm:moment-control}
Let
\[
H_0=\frac12\sum_{j=1}^n\omega_j(p_j^2+q_j^2-1),\qquad \omega_j>0,
\]
on $L^2(\mathbb R^n)$, and let $P$ be a real polynomial of even degree
$N$ satisfying $P(x)\ge\varepsilon|x|^N-k$ for some
$\varepsilon,k>0$.  There is $\Lambda>0$ such that, for every
$0<\lambda_0\le\Lambda$, exactly one normalized nonnegative
Schwartz-valued norm-continuous map
$u:[0,\lambda_0]\to L^2(\mathbb R^n)$ starts at the free vacuum and has
$u(\lambda)$ as the ground state of the self-adjoint closure of
\[
H_0+\lambda:P(q):_{u(\lambda)}.
\]
In particular, no a priori bound on
$\langle u(\lambda),P(q)u(\lambda)\rangle$ is required.
\end{theorem}

\section{State-relative Wick polynomials as a formal Appell family}
\label{sec:appell}

For a normalized Schwartz vector $u$ write
\begin{equation}
  \label{eq:moments}
  \mu_\alpha(u)=\langle u,q^\alpha u\rangle,
  \qquad
  A_\alpha^u(q)=\,:q^\alpha:_u .
\end{equation}
The defining expectation and commutator axioms for state-relative
renormalized products \cite{BaezZhou1992} are
\begin{equation}
  \label{eq:axioms}
  A_0^u=1,
  \qquad
  \langle u,A_\alpha^u(q)u\rangle=0\ \ (|\alpha|>0),
  \qquad
  \partial_{q_j}A_\alpha^u=\alpha_jA_{\alpha-e_j}^u .
\end{equation}

\begin{proposition}
  \label{prop:generating}
  Let
  \begin{equation}
    \label{eq:series}
    \mathcal A_u(t,q)=\sum_{\alpha\ge0}A_\alpha^u(q)\frac{t^\alpha}{\alpha!},
    \qquad
    M_u(t)=\sum_{\alpha\ge0}\mu_\alpha(u)\frac{t^\alpha}{\alpha!}
  \end{equation}
  as formal power series. Then
  \begin{equation}
    \label{eq:appell}
    \mathcal A_u(t,q)=\frac{e^{t\cdot q}}{M_u(t)} .
  \end{equation}
\end{proposition}

\begin{proof}
  The third axiom in \eqref{eq:axioms} gives
  $\partial_{q_j}\mathcal A_u=t_j\mathcal A_u$, so
  $\mathcal A_u(t,q)=e^{t\cdot q}C_u(t)$ for a series $C_u$ not involving
  $q$. Taking the $u$-expectation of \eqref{eq:series} coefficientwise and
  using the second axiom gives $M_u(t)C_u(t)=1$; since $M_u(0)=1$, the
  formal reciprocal exists and is unique.
\end{proof}

No analytic moment-generating function is asserted or needed: every
identity here is an identity of formal series, and each coefficient is a
finite algebraic expression in finitely many moments.

\begin{proposition}[Change of reference state]
  \label{prop:change}
  Let $v$ be another normalized Schwartz vector and put
  $r_\alpha(v,u)=\langle v,A_\alpha^u(q)v\rangle$. Then
  \begin{equation}
    \label{eq:change}
    \sum_\alpha r_\alpha(v,u)\frac{t^\alpha}{\alpha!}
    =\frac{M_v(t)}{M_u(t)},
  \end{equation}
  equivalently
  \begin{equation}
    \label{eq:recurrence}
    r_\alpha(v,u)=\mu_\alpha(v)-\mu_\alpha(u)
      -\sum_{\substack{0<\beta\le\alpha\\ \beta\ne\alpha}}
        \binom{\alpha}{\beta}\mu_\beta(u)\,r_{\alpha-\beta}(v,u).
  \end{equation}
\end{proposition}

\begin{proof}
  Take the $v$-expectation of \eqref{eq:appell} coefficientwise, which gives
  \eqref{eq:change}; equating coefficients in
  $M_u(t)\sum_\alpha r_\alpha t^\alpha/\alpha!=M_v(t)$ gives
  \eqref{eq:recurrence}.
\end{proof}

\section{Concentration under top-moment escape}
\label{sec:concentration}

\begin{lemma}
  \label{lem:measure}
  Let probability measures $\nu_j\to\nu$ in total variation and suppose
  $M_j=\int|x|^N\,d\nu_j\to\infty$. Then for every $0<r<N$,
  \begin{equation}
    \label{eq:lower-moments}
    \int|x|^r\,d\nu_j=o\bigl(M_j^{r/N}\bigr).
  \end{equation}
\end{lemma}

\begin{proof}
  Put $\theta=r/N$. For any $R>0$, Hölder's inequality on the tail gives
  \begin{equation*}
    \int|x|^r\,d\nu_j\le R^r+M_j^{\theta}\,\nu_j(|x|>R)^{1-\theta}.
  \end{equation*}
  Dividing by $M_j^{\theta}$ and taking the upper limit, total-variation
  convergence gives
  \begin{equation*}
    \limsup_{j\to\infty}\frac{\int|x|^r\,d\nu_j}{M_j^{r/N}}
    \le\nu(|x|>R)^{1-r/N}.
  \end{equation*}
  Letting $R\to\infty$ makes the right side vanish.
\end{proof}

\begin{lemma}
  \label{lem:tv}
  Let normalized $u_j\to v$ in $L^2(\mathbb R^n)$ and suppose
  $M_j=\langle u_j,|q|^Nu_j\rangle\to\infty$. Then the position laws of
  $u_j$ converge to that of $v$ in total variation, and
  \begin{equation}
    \label{eq:mu-small}
    \mu_\alpha(u_j)=o\bigl(M_j^{|\alpha|/N}\bigr)
    \qquad(0<|\alpha|<N).
  \end{equation}
\end{lemma}

\begin{proof}
  The Cauchy--Schwarz inequality gives
  \begin{equation}
    \label{eq:tv-bound}
    \int\bigl||u_j|^2-|v|^2\bigr|
    \le\lVert u_j-v\rVert_2\,\lVert u_j+v\rVert_2
    \le2\lVert u_j-v\rVert_2\longrightarrow0,
  \end{equation}
  which is total-variation convergence of the position laws.
  \Cref{lem:measure} together with $|x^\alpha|\le|x|^{|\alpha|}$ then gives
  \eqref{eq:mu-small}.
\end{proof}

\section{A coercive negative change-of-state asymptotic}
\label{sec:coercive}

\begin{lemma}
  \label{lem:r-asymptotics}
  Fix a normalized Schwartz vector $v$ and let normalized Schwartz vectors
  $u_j\to v$ in norm satisfy $M_j\to\infty$. Writing
  $r_\alpha^{(j)}=r_\alpha(v,u_j)$,
  \begin{equation}
    \label{eq:r-small}
    r_\alpha^{(j)}=o\bigl(M_j^{|\alpha|/N}\bigr)
    \quad(0<|\alpha|<N),
    \qquad
    r_\alpha^{(j)}=-\mu_\alpha(u_j)+o(M_j)
    \quad(|\alpha|=N).
  \end{equation}
\end{lemma}

\begin{proof}
  The moments of $v$ are fixed, so for every $0<|\gamma|<N$ we also have
  $\mu_\gamma(v)=o(M_j^{|\gamma|/N})$. Induct on $|\alpha|$ in
  \eqref{eq:recurrence} using \eqref{eq:mu-small}: each term of the sum is a
  product of two quantities whose little-$o$ weights add to $|\alpha|$, which
  gives the first assertion. At total degree $N$ the same recurrence leaves
  $-\mu_\alpha(u_j)$ as the only contribution of top weight, all other terms
  being $o(M_j)$.
\end{proof}

\begin{proposition}
  \label{prop:to-minus-infinity}
  Under the hypotheses of \Cref{lem:r-asymptotics},
  \begin{equation}
    \label{eq:blowdown}
    \langle v,:P(q):_{u_j}v\rangle\ \le\ -\varepsilon M_j+o(M_j)
    \ \longrightarrow\ -\infty .
  \end{equation}
\end{proposition}

\begin{proof}
  Write $P=P_N+P_{<N}$ with $P_N$ homogeneous of degree $N$. By linearity of
  the renormalization, \eqref{eq:r-small} gives
  \begin{equation}
    \label{eq:leading}
    \langle v,:P(q):_{u_j}v\rangle
    =-\langle u_j,P_N(q)u_j\rangle+o(M_j).
  \end{equation}
  Evaluating the coercivity hypothesis $P(x)\ge\varepsilon|x|^N-k$ at $tx$,
  dividing by $t^N$ and letting $t\to\infty$ gives
  $P_N(x)\ge\varepsilon|x|^N$, so
  $\langle u_j,P_N(q)u_j\rangle\ge\varepsilon M_j$.
\end{proof}

\section{The opposing ground-state bound}
\label{sec:opposing}

\begin{lemma}
  \label{lem:lower}
  Let $u$ be a normalized Schwartz ground state of
  $H_u=H_0+\lambda:P(q):_u$ with $\lambda>0$. Then for every normalized
  Schwartz vector $v$,
  \begin{equation}
    \label{eq:lower-bound}
    \langle v,:P(q):_u v\rangle
    \ \ge\ P(0)-\frac{\langle v,H_0v\rangle}{\lambda},
  \end{equation}
  and for the free vacuum $u_0$, uniformly in $\lambda$,
  \begin{equation}
    \label{eq:vacuum-bound}
    \langle u_0,:P(q):_u u_0\rangle\ \ge\ P(0).
  \end{equation}
\end{lemma}

\begin{proof}
  The expectation normalization in \eqref{eq:axioms}, including
  $:1:_u=1$, gives $\langle u,:P(q):_u u\rangle=P(0)$, so the ground energy
  satisfies
  $E_u=\langle u,H_0u\rangle+\lambda P(0)\ge\lambda P(0)$ because
  $H_0\ge0$. Ground-state minimality against $v$ gives
  $E_u\le\langle v,H_0v\rangle+\lambda\langle v,:P(q):_u v\rangle$, and
  combining the two yields \eqref{eq:lower-bound}. Taking $v=u_0$ and using
  $\langle u_0,H_0u_0\rangle=0$ gives \eqref{eq:vacuum-bound}, in which
  $\lambda$ no longer appears.
\end{proof}

The two preceding sections now point in opposite directions, and that is the
whole mechanism: \Cref{prop:to-minus-infinity} sends the same expectation to
$-\infty$ whenever the top moment escapes, while \Cref{lem:lower} holds it
above a fixed number.

\section{Automatic boundedness of the omitted moment}
\label{sec:automatic}

\begin{proposition}
  \label{prop:bounded}
  Let $u:[0,L]\to L^2(\mathbb R^n)$ be any normalized nonnegative
  Schwartz-valued norm-continuous self-consistent branch on a compact
  coupling interval. Then
  \begin{equation}
    \label{eq:top-bounded}
    \sup_{0\le\lambda\le L}\langle u(\lambda),|q|^Nu(\lambda)\rangle<\infty,
  \end{equation}
  and consequently
  \begin{equation}
    \label{eq:P-bounded}
    \sup_{0\le\lambda\le L}
      \bigl|\langle u(\lambda),P(q)u(\lambda)\rangle\bigr|<\infty .
  \end{equation}
\end{proposition}

\begin{proof}
  Suppose the radial top moment were unbounded, so that there are
  $\lambda_j\in[0,L]$ with
  $M_j=\langle u(\lambda_j),|q|^Nu(\lambda_j)\rangle\to\infty$. Passing to a
  subsequence, $\lambda_j\to\lambda_*$ by compactness, and norm continuity
  gives $u_j:=u(\lambda_j)\to v:=u(\lambda_*)$ in $L^2$. The values
  $\lambda_j=0$ cannot witness escape, since $u(0)$ is a fixed vector, so we
  may take $\lambda_j>0$. By \Cref{prop:to-minus-infinity},
  \begin{equation}
    \label{eq:Wj}
    W_j:=\langle v,:P(q):_{u_j}v\rangle\longrightarrow-\infty .
  \end{equation}

  If $\lambda_*>0$ then eventually $\lambda_j\ge\lambda_*/2$, and
  \eqref{eq:lower-bound} with the fixed test vector $v$ gives
  $W_j\ge P(0)-2\langle v,H_0v\rangle/\lambda_*$, contradicting
  \eqref{eq:Wj}. If $\lambda_*=0$ then continuity and uniqueness of the
  normalized nonnegative free ground state give $v=u(0)=u_0$, and
  \eqref{eq:vacuum-bound} gives $W_j\ge P(0)$, again contradicting
  \eqref{eq:Wj}. This proves \eqref{eq:top-bounded}, and
  \eqref{eq:P-bounded} follows from $|P(x)|\le C_P(1+|x|^N)$.
\end{proof}

\section{Uniform uniqueness on every sufficiently short interval}
\label{sec:uniqueness}

We use the finite-dimensional cone inverse construction of Baez and Zhou
\cite[Theorem 1 and Lemmas 8--9]{BaezZhou1992}, keeping track
of the interval quantifier.

Let $\mathcal P_N$ be the finite-dimensional space of real polynomials of
degree at most $N$, and let $\mathcal C_+\subset\mathcal P_N$ consist of
those whose degree-$N$ homogeneous form is strictly positive on the unit
sphere; this is an open convex cone contained in the cone of polynomials
bounded below. For $Q\in\mathcal C_+$ let $v(Q)$ be the normalized
nonnegative ground state of $H_0+Q(q)$ and define $T(Q)$ by
\begin{equation}
  \label{eq:T}
  Q(q)=\,:T(Q)(q):_{v(Q)} .
\end{equation}
The cited source proves that $T$ is analytic on $\mathcal C_+$ and that
$T(Q)\to0$ and $dT(Q)\to I$ as $Q\to0$ within the cone, and its cone inverse
lemma supplies $\Lambda>0$ and a continuous curve
\begin{equation}
  \label{eq:curve}
  y:[0,\Lambda]\to\mathcal C_+\cup\{0\},
  \qquad
  T(y(\lambda))=\lambda P,
  \qquad
  y(0)=0 .
\end{equation}
That lemma's uniqueness argument applies verbatim on every restricted
interval $[0,L]$ with $0<L\le\Lambda$: near zero the contraction estimate
forces any continuous solution of \eqref{eq:curve} to equal $y$, and if
equality first failed later, local invertibility of $T$ at the common
first-disagreement point would give a contradiction. So $y|_{[0,L]}$ is the
unique continuous cone-valued solution on each such interval.

\begin{lemma}
  \label{lem:Qw-continuous}
  Fix $0<L\le\Lambda$ and let $w:[0,L]\to L^2(\mathbb R^n)$ be any branch
  allowed in \Cref{thm:moment-control}. Set
  \begin{equation}
    \label{eq:Qw}
    Q_w(0)=0,
    \qquad
    Q_w(\lambda)=\lambda:P(q):_{w(\lambda)}\quad(\lambda>0).
  \end{equation}
  Then $Q_w$ takes values in $\mathcal C_+\cup\{0\}$, is continuous
  coefficientwise on $[0,L]$ including at $\lambda=0$, and satisfies
  $T(Q_w(\lambda))=\lambda P$.
\end{lemma}

\begin{proof}
  The leading form of $Q_w(\lambda)$ is $\lambda P_N$, so
  $Q_w(\lambda)\in\mathcal C_+$ for $\lambda>0$, and self-consistency with
  \eqref{eq:T} gives $w(\lambda)=v(Q_w(\lambda))$ and
  $T(Q_w(\lambda))=\lambda P$.

  For continuity, write
  $B=\sup_{0\le\lambda\le L}\langle w(\lambda),|q|^Nw(\lambda)\rangle$, which
  is finite by \Cref{prop:bounded}. Norm continuity gives total-variation
  continuity of the position laws by \eqref{eq:tv-bound}, and for
  $0\le r<N$ the uniform tail estimate
  \begin{equation}
    \label{eq:tail}
    \int_{|x|>R}|x|^r|w(\lambda,x)|^2\,dx\ \le\ R^{r-N}B
  \end{equation}
  holds. Truncating at $R$ and letting $R\to\infty$ shows every moment of
  total degree below $N$ is continuous in $\lambda$. By
  \eqref{eq:recurrence}, every nonconstant coefficient of
  $:P:_{w(\lambda)}$ depends only on moments of total degree at most $N-1$,
  so the nonconstant part $Q_w^{\mathrm{nc}}$ is continuous, including at
  zero.

  For the constant coefficient $c_w$, adding a scalar to a potential changes
  neither its ground state nor any positive-degree renormalized product, so
  $T(R+c)=T(R)+c$. With $T(Q_w(\lambda))=\lambda P$ this gives
  \begin{equation}
    \label{eq:cw}
    c_w(\lambda)=\lambda P(0)-\bigl[T(Q_w^{\mathrm{nc}}(\lambda))\bigr]_0,
  \end{equation}
  where $[\,\cdot\,]_0$ is the constant coefficient. Continuity of $T$ on the
  cone, together with $T(Q)\to0$ as $Q\to0$, makes $c_w$ continuous on
  $[0,L]$.
\end{proof}

\begin{proof}[Proof of \Cref{thm:moment-control}]
  Let $\Lambda$ be as in \eqref{eq:curve} and fix $0<\lambda_0\le\Lambda$.
  Existence of a branch on $[0,\lambda_0]$ follows by restricting
  $y$ and setting $u(\lambda)=v(y(\lambda))$.

  For uniqueness, let $w$ be any branch allowed by the statement. By
  \Cref{lem:Qw-continuous}, $Q_w$ is a continuous cone-valued solution of
  \eqref{eq:curve} on $[0,\lambda_0]$, so restricted-interval cone
  uniqueness gives $Q_w=y|_{[0,\lambda_0]}$ and hence
  $w(\lambda)=v(y(\lambda))$ for every $\lambda$.

  No bound on $\langle w(\lambda),P(q)w(\lambda)\rangle$ was assumed at any
  point: the one used in \Cref{lem:Qw-continuous} was supplied by
  \Cref{prop:bounded} from norm continuity and self-consistency alone.
\end{proof}

\section{Conclusion}

Norm continuity, coercivity, and ground-state minimality together prevent
escape of the interaction moment.  The conclusion is local in coupling and
finite dimensional; it neither supplies maximal continuation nor removes
continuity at the free-vacuum endpoint.  Within that boundary, the formerly
assumed moment condition is redundant.

\section*{Acknowledgments}

Z.\ Y.\ acknowledges support from the Strategic Priority Research Program of the
Chinese Academy of Sciences under Grant No.~XDB1680102. Z.\ Y.\ also
acknowledges support by National Natural Science Foundation of China (Grant
No.~12247101). We also acknowledge Gewu Intelligence Lab for providing the
collaborative research environment in which this project was developed.

OpenAI GPT-5.6 Sol was used substantively in this work: in solving the
problem, including the formulation of arguments and the derivation of
intermediate steps, and in drafting and revising the manuscript. The
authors specified the problem, guided the approach to be taken, checked
every argument and result it produced, and take full responsibility for
the content of this manuscript.

% Numbers the reference list in order of first citation, so [1] is the first
% work the text cites.
\bibliographystyle{unsrt}
\bibliography{references,unverified}
\end{document}
